\chapter{Formal setup, notation, and model specification}
\label{ch:formal}
\chapterpromise{Ordered segmentation, exposure-aware local likelihoods, posterior targets, and the exact factorization required by the partition dynamic program.}

\section{Notation and symbols}
\label{sec:notation}

This section fixes notation used throughout the book. We use lower-case Roman letters for
indices, reserve $t=(t_0,\dots,t_k)$ for boundary vectors, and adopt the half-open block notation
$(i,j]:=\{i+1,\dots,j\}$ for $0\le i<j\le n$.

For irregular designs, physical segment lengths are defined using candidate boundary coordinates
$u_0<u_1<\cdots<u_n$ rather than by implicitly referencing an undefined observation location
$x_0$. The physical length of block $(i,j]$ is
$\Delta_x(i,j):=u_j-u_i$. On a regular index grid set $u_i=i$; for interval data the
$u_i$ are the observed interval endpoints. Point observations without meaningful physical interval
endpoints use the index-uniform prior or provide an external left endpoint $u_0$.

\paragraph{Conjugate-prior normalization convention.}
For exponential-family blocks we use the following (Diaconis--Ylvisaker) convention for the
conjugate prior normalizer. For hyperparameters $(\alpha,\beta)$,
\begin{equation}
\begin{aligned}
p(\theta\mid \alpha,\beta)
  &= \exp\!\left\{\eta(\theta)^\top \alpha-\beta A(\theta)-\log Z(\alpha,\beta)\right\},\\
Z(\alpha,\beta)
  &:= \int \exp\!\left\{\eta(\theta)^\top \alpha-\beta A(\theta)\right\}\,d\theta.
\end{aligned}
\label{eq:Zconvention}
\end{equation}
With this convention, $Z$ is the integral itself (not its reciprocal) and block evidences take
the ratio form $Z(\alpha_0+S_{ij},\beta_0+W_{ij})/Z(\alpha_0,\beta_0)$ (Theorem~\ref{thm:ef-integral}
and Chapter~\ref{ch:block}).

\begin{longtable}{@{}p{0.22\linewidth}p{0.72\linewidth}@{}}
\caption[Notation and symbols]{Notation and symbols used throughout the book. This is a notation glossary; no source CSV is associated with it.}\label{tab:notation}\\
\toprule
\textbf{Symbol} & \textbf{Meaning} \\
\midrule
$n$ & Number of observations on the ordered design. \\
$x_i$ & Design point / location of observation $y_i$. \\
$u_i$ & Candidate boundary coordinate used for physical segment lengths, with $u_0<\cdots<u_n$. \\
$\Delta_x(i,j)$ & Physical length of block $(i,j]$, defined as $u_j-u_i$. \\
$y_i$ & Observation at location $x_i$ (single-sequence case). \\
$d_i$ & Family-specific exposure, trial-count, offset, or precision descriptor when required by the observation model. Optional power/reliability weights are denoted $a_i$; neither role is inferred from irregular spacing. \\
\midrule
$k$; $k_{\max}$ & Number of segments; maximum segment count considered by the DP. \\
$t=(t_0,\dots,t_k)$ & Boundary vector, with $0=t_0<t_1<\cdots<t_k=n$. \\
$I_q(t)$ & Index set of segment $q$: $\{t_{q-1}+1,\dots,t_q\}$. \\
$(i,j]$ & Candidate block $\{i+1,\dots,j\}$. \\
\midrule
$S(y,d)$, $E(d)$, $\eta(\theta)$, $A(\theta)$, $h(y,d)$ & Family-specific statistic contribution, exposure multiplier, natural-parameter map, log-partition function, and base measure for the chosen exponential-family block model. \\
$m_\star(\theta)=m(\theta;d_\star)$ & Declared segment functional on the reporting scale, evaluated at a fixed reference descriptor $d_\star$ when the observation mean depends on exposure or trial structure. Examples include a Gaussian mean, a Poisson rate at unit exposure, and a Binomial success probability. \\
\midrule
$(\alpha,\beta)$ & Conjugate hyperparameters with normalizer $Z(\alpha,\beta)$ defined in \eqref{eq:Zconvention}. \\
$(\alpha_0,\beta_0)$ & Prior hyperparameters before observing block data. \\
\midrule
$S_{ij}$ & Block statistic sum $\sum_{t=i+1}^j a_t S(y_t,d_t)$. \\
$W_{ij}$ & Block log-partition multiplier $\sum_{t=i+1}^j a_t E(d_t)$. \\
$H_{ij}$ & Powered base-measure sum $\sum_{t=i+1}^j a_t\log h(y_t,d_t)$. \\
$A^{(0)}_{ij}$ & Block evidence on $(i,j]$. \\
$A^{(r)}_{ij}$ & Block moment numerator $A^{(0)}_{ij}\,\mathbb{E}[m_\star(\theta)^r\mid y_{i+1:j}]$ (for integer $r\in\{1,2,\dots\}$; conjugate-case definition). In the non-conjugate Laplace discussion of Section~\ref{sec:nonconj}, the superscript is a test function rather than a moment order; those objects are denoted $M^{[g]}_{ij}$ to avoid collision. \\
$\widetilde A^{(r)}_{ij}$ & Prior-adjusted block quantity $A^{(r)}_{ij}\gamma_x(i,j)$, where $\gamma_x$ combines cohesion and an interior-boundary hazard when appropriate. \\
\midrule
$L_{k,j}$ & Prefix evidence for $y_{1:j}$ with exactly $k$ segments. \\
$R_{k,i}$ & Suffix evidence for $y_{i+1:n}$ with exactly $k$ segments. \\
$C_k$ & Normalizer of $p(t\mid k)$ under a segment-factorized partition prior. \\
$p(k\mid y)$ & Posterior over segment counts. \\
$p(t_p=h\mid y,k)$ & Marginal posterior of the $p$th boundary given $k$. \\
$p(b_i=1\mid y,k)$ & Boundary-event marginal at interior index $i$, equal to $\sum_{p=1}^{k-1} p(t_p=i\mid y,k)$. This is the calibration target in Section~\ref{sec:experiments}. \\
\midrule
$\widehat t^{\mathrm{MAP}}(k)$ & Joint MAP segmentation given $k$, obtained by max-sum DP and backtracking. \\
$\widehat\mu_q^{(r)}$ & Posterior moment of the segment-$q$ mean parameter conditional on a chosen segmentation. \\
$\widehat f_t^{\mathrm{Bayes}}$ & Bayes regression curve at index $t$ (posterior mean of the latent piecewise-constant signal). \\
\midrule
$s\in\{1,\dots,S\}$ & Sequence / subject index in multi-sequence settings. \\
$g\in\{1,\dots,G\}$ & Group index. \\
$r_{sg}$ & Auxiliary allocation weight for sequence $s$ and latent group $g$ in the finite latent-group criterion. \\
\midrule
$\mathcal{M}_g$ & Fitted group-specific model exported for prediction. \\
$\ell_g$ & Log predictive score of new data under group $g$. \\
$\pi_g$ & Prior group weight used for prediction or in the score-template objective. \\
\midrule
$\varepsilon$ & Uniform per-block log-evidence error tolerance for non-conjugate block approximations; appears as the hypothesis controlling the posterior-odds stability bound in Section~\ref{sec:nonconj}. \\
$\widehat A^{(0)}_{ij}$ & Approximate (non-conjugate) block evidence; satisfies $|\log\widehat A^{(0)}_{ij} - \log A^{(0)}_{ij}| \le \varepsilon$ under the stability assumption. \\
$\phi_t$ & Known per-observation precision used in the Beta-block likelihood for $y_t\in(0,1)$. \\
$\omega_t$ & P\'olya--Gamma augmentation variable for logistic / negative-binomial mean-field block approximations. \\
$M^{[g]}_{ij}$ & Block test-function numerator $\int g(\theta)\, p(y_{i+1:j}\mid \theta)\, p(\theta)\, d\theta$, used in the Laplace-approximation discussion of Section~\ref{sec:nonconj} where the superscript denotes a test function rather than a moment order. \\
\midrule
$f_{\mathrm{Lap}}$, $f_{\mathrm{JJ}}$, $f_{\mathrm{EP}}$, $f_{\mathrm{PG}}$, $f_{\mathrm{Q}}$ & Local block-evidence approximations:
Laplace, Jaakkola--Jordan variational bound, expectation propagation, P\'olya--Gamma mean field,
and one-dimensional quadrature. Each instantiates $\widehat A^{(0)}_{ij}$ and inherits the same DP layer. \\
\bottomrule
\end{longtable}

\paragraph{Remarks.}
Two recurring distinctions: \emph{design irregularity} may enter through the partition prior
(\S\ref{sec:irregular}); $d_i$ enters the likelihood only when the observation model supplies
genuine exposures or known precisions. \emph{Boundary marginals} and the \emph{joint MAP
segmentation} are different posterior summaries (\S\ref{sec:dp}) and must not be conflated.

\label{sec:problem}

BayesBreak addresses \emph{offline} Bayesian segmentation of ordered data. Throughout, the
fundamental unknown is a contiguous partition of the index set $\{1,\dots,n\}$ together with a
segment parameter attached to each block. This section states the objects of interest in a
model-agnostic way; Section~\ref{sec:notation} introduces the detailed notation used later.

Unless stated otherwise, all inferential statements in the book use the following standing assumption.

\begin{assumption}[Offline segmentation support and block independence]
\label{ass:standing-offline}
The observations are ordered on a fixed set of candidate boundary locations and the segment-count
set $\{1,\dots,k_{\max}\}$ is finite. Structural admissibility of a candidate block is determined
before observing the response values by the candidate grid, declared segment-length restrictions,
and prior support. Every structurally admissible block has finite nonnegative block evidence; a zero
evidence contributes zero to the likelihood-weighted partition sum but does not change the prior
support or its normalizer $C_k$. Structurally excluded blocks have local prior factor zero and are
stored with log-score $-\infty$. Segment parameters are conditionally independent across blocks
given a segmentation, and the conditional partition prior factorizes over segments as in
Assumption~\ref{ass:factorizing-prior}. For approximate non-conjugate blocks, the same statement
applies to the supplied approximate score array, while comparison with the exact non-conjugate model
requires the additional error hypothesis of Assumption~\ref{ass:uniform-block-error}.
\end{assumption}

\subsection{Single-sequence segmentation and Bayes regression}
\label{sec:problem-single}

We observe ordered pairs $\{(x_i,y_i)\}_{i=1}^n$ with strictly increasing design points
$x_1<\cdots<x_n$. The design can be irregular. When the observation model includes known
exposures, trial counts, or precision weights, we also allow deterministic nonnegative quantities
$w_i$; unless explicitly stated otherwise, point-observation models use $w_i\equiv 1$.

A segmentation into $k$ contiguous blocks is represented by a boundary vector
\begin{equation}
0=t_0<t_1<\cdots<t_{k-1}<t_k=n,
\end{equation}
where block $q\in\{1,\dots,k\}$ is the index set
\begin{equation}
I_q(t)\coloneqq \{t_{q-1}+1,\dots,t_q\}.
\end{equation}

\begin{definition}[Segmentation space and structural support]
\label{def:segmentation-space}
For fixed $n$ and $k$, let
\[
\mathcal{T}_{k,n}:=\{t=(t_0,\dots,t_k):0=t_0<t_1<\cdots<t_k=n\}.
\]
The \emph{structural support} $\mathcal{T}^{+}_{k,n}\subseteq\mathcal{T}_{k,n}$ consists of
those boundary vectors for which every block satisfies the declared grid and length restrictions
and has positive local prior factor $\gamma_x(t_{q-1},t_q)$. Conditional on the design and model
specification, this set is fixed independently of the observed response values and of
$A^{(0)}_{ij}$. A structurally admissible partition may have zero likelihood for the observed data;
it then contributes zero to the posterior partition sum, while the prior normalizer $C_k$ continues
to sum over the full structural support.
\end{definition}

The corresponding segment parameter is denoted $\theta_q\in\Theta$. Conditionally on
$(k,t,\theta_{1:k})$, observations factorize across blocks:
\begin{equation}
p(y_{1:n}\mid k,t,\theta_{1:k})
= \prod_{q=1}^k \prod_{i\in I_q(t)} p(y_i\mid \theta_q, w_i),
\label{eq:problem-lik}
\end{equation}
where $p(y_i\mid\theta_q,w_i)$ is interpreted according to the chosen family. For example,
depending on the family, $w_i$ encodes exposure in Poisson models, trial count in Binomial models, or known precision
in Gaussian models. Irregular spacing by itself does not force a non-unit $w_i$; design geometry
can instead enter through the partition prior.

For any candidate block $(i,j]$ with $0\le i<j\le n$, the key primitive is the integrated
single-block evidence
\begin{equation}
\mathcal{L}(i,j;\alpha_0,\beta_0)
\coloneqq p(y_{i+1:j}\mid \alpha_0,\beta_0)
= \int \prod_{t=i+1}^j p(y_t\mid\theta,w_t)\,p(\theta\mid\alpha_0,\beta_0)\,d\theta,
\label{eq:problem-block-evidence}
\end{equation}
where $p(\theta\mid\alpha_0,\beta_0)$ is the segment prior with hyperparameters $(\alpha_0,\beta_0)$
(see Section~\ref{sec:notation}). When $\mathcal{L}(i,j)$ is available in closed form, the
entire offline posterior can be computed exactly by dynamic programming. When it is only
approximated, BayesBreak isolates the approximation at the block level and reuses the same global
DP.

\begin{definition}[Boundary-event variables, Bayes curve, and exported estimator]
\label{def:problem-bayes-curve}
For a boundary vector $t\in\mathcal{T}_{k,n}$, define the interior boundary indicator
$b_i(t):=\sum_{p=1}^{k-1}\mathbf{1}\{t_p=i\}$ for $i=1,\dots,n-1$. Because the boundaries are
strictly ordered, $b_i(t)\in\{0,1\}$. For fixed $k$, the Bayes regression curve at index $i$ is
the posterior mean of the observation-scale segment mean covering $i$ after marginalizing over
$t\in\mathcal{T}^{+}_{k,n}$ and over segment parameters. The exported segmentation estimator is a
deterministic boundary vector, usually the joint MAP vector $\widehat t^{\mathrm{MAP}}(k)$, paired
with segment-level posterior summaries for downstream prediction. These two objects need not
coincide because the Bayes curve averages over boundary uncertainty whereas the exported estimator
chooses one segmentation.
\end{definition}

The inferential targets are:
\begin{enumerate}[label=(\roman*), leftmargin=2.3em]
    \item the marginal evidence at fixed segment count,
    \[
      p(y_{1:n}\mid k)
      =\sum_{t\in\mathcal{T}^{+}_{k,n}}p(t\mid k)
        \prod_{q=1}^{k}\mathcal{L}(t_{q-1},t_q);
    \]
    \item the posterior $p(k,t\mid y)$ over segment counts and boundaries and its marginals
    $p(k\mid y)$, $p(t_p=h\mid y,k)$, and the boundary-event marginal
    \(p(b_i=1\mid y,k)=\sum_{p=1}^{k-1}p(t_p=i\mid y,k)\)
    (the calibration target of \S\ref{sec:experiments});
    \item posterior summaries of the latent piecewise-constant signal, namely
    $p(\theta_i\mid y)$ and $\mathbb{E}[m_\star(\theta_i)\mid y]$ with $m_\star$ the
    observation-scale mean parameter (the Bayes regression curve; conditional on the selected
    $\widehat k$ unless stated otherwise);
    \item point estimators for downstream use. In particular the \emph{joint} MAP segmentation
    \(\widehat t^{\mathrm{MAP}}(k)\in\arg\max_t p(t\mid y,k)\)
    differs from the vector of marginal boundary modes
    \(\arg\max_h p(t_p=h\mid y,k)\)
    and is recovered by max-sum DP. The prior on $k$ and the normalizer $\log C_k$ enter the
    objective only when also optimizing over $k$.
\end{enumerate}

The three posterior summaries entering downstream analysis are isolated as labelled objects so
that later sections can refer to them without re-deriving notation.

\begin{definition}[Posterior summaries: marginal boundaries, segment count, and joint MAP]
\label{def:posterior-summaries}
Fix a sequence $y=y_{1:n}$, a maximum segment count $k_{\max}$, and Assumption~\ref{ass:standing-offline}.
Define:
\begin{enumerate}[label=(\alph*), leftmargin=2.3em]
    \item the \emph{marginal boundary posterior}
    \(\mathrm{MB}_i(k):=p(b_i=1\mid y,k)\)
    for interior indices $i\in\{1,\dots,n-1\}$ and admissible $k\in\{1,\dots,k_{\max}\}$;
    \item the \emph{posterior over segment count}
    \(\mathrm{PK}(k):=p(k\mid y)\propto p(k)\,p(y\mid k)\)
    for $k\in\{1,\dots,k_{\max}\}$;
    \item the \emph{joint MAP segmentation at fixed $k$}
    \(\widehat t^{\mathrm{MAP}}(k):=\arg\max_{t\in\mathcal{T}^{+}_{k,n}} p(t\mid y,k)\),
    with ties broken by a deterministic rule (Section~\ref{sec:dp-alg}).
\end{enumerate}
These three objects are returned by distinct DP recursions (sum-product, sum-product with
$k$-aggregation, and max-sum, respectively), and they are not in general consistent with each
other: $\mathrm{MB}_i$ and $\widehat t^{\mathrm{MAP}}(k)$ can place boundary mass at different
indices (Remark~\ref{rem:marg-vs-joint}), and the mode of $\mathrm{PK}$ need not coincide with
the cardinality of any single thresholded set of marginal boundaries.
\end{definition}

\subsection{Multi-sequence structure}
\label{sec:problem-multiseq}

Many applications provide multiple related sequences observed on a common ordered grid or on a
common set of candidate boundary locations. Let sequence $s\in\{1,\dots,S\}$ be denoted by
$y^{(s)}_{1:n}$, with optional weights $w_i^{(s)}$ and subject-specific segment parameters
$\theta_q^{(s)}$ when boundaries are shared. We use $z_s\in\{1,\dots,G\}$ for group membership,
$\mathcal{S}_g:=\{s:z_s=g\}$ for known groups, and $r_{sg}$ for latent-group responsibilities.
BayesBreak treats three progressively richer settings.

\begin{definition}[Multi-sequence segmentation classes]
\label{def:multiseq-classes}
A shared-boundary replicate model has one boundary vector $t$ and subject-specific segment
parameters $\theta^{(s)}_{1:k}$. A known-group model has observed labels $z_s$ and one boundary
vector $t^{(g)}$ per group. A latent-template model has unobserved labels $z_s$ and a finite set of
templates $\tau_g=(k_g,t^{(g)})$ optimized under the finite latent-group criterion of
Section~\ref{sec:latent-em}. These three classes preserve the block-evidence interface; fully
Bayesian averaging over both unknown labels and all group segmentations is a different target and
is not claimed for the score-template model.
\end{definition}

\paragraph{Shared-boundary replicates.}
All sequences share a common segmentation $t\in\mathcal{T}^{+}_{k,n}$, but each sequence has its
own within-block parameters $\theta^{(s)}_{1:k}$. The common-grid requirement is an assumption on
the candidate boundary coordinates; subject-specific missingness or exposures are handled inside
the block evidences. Conditional on $t$, the sequences are independent:
\begin{equation}
p(\{y^{(s)}\}_{s=1}^S\mid t,\{\theta^{(s)}\})
= \prod_{s=1}^S \prod_{q=1}^k \prod_{i\in I_q(t)} p\bigl(y_i^{(s)}\mid \theta_q^{(s)}, w_i^{(s)}\bigr).
\end{equation}
Integrating out the subject-specific segment parameters yields pooled block evidences obtained by
multiplying subject-level block evidences.

\paragraph{Known groups.}
Sequences are partitioned into known groups $g\in\{1,\dots,G\}$. Each group has its own
segmentation $t^{(g)}\in\mathcal{T}^{+}_{k_g,n}$, while sequences within a group share that
segmentation and retain their own segment parameters. Conditional independence is within group
and across subjects given $t^{(g)}$.

\paragraph{Latent groups.}
Group labels are unknown. A fully Bayesian formulation averages over both
assignments and group-specific segmentations. BayesBreak adopts a tractable
\emph{template-mixture} formulation: each group is associated with a single latent boundary
template $\tau_g=(k_g,t^{(g)})$, and coordinate ascent alternates between soft assignments (normalized auxiliary weights for the current
templates) and exact responsibility-weighted max-sum DP updates of those templates. The
finite latent-group criterion is non-decreasing under the exact alternating updates
(Theorem~\ref{thm:em-monotone}); this is a finite-template optimization model rather than exact Bayesian
averaging over latent-group segmentations.

\subsection{Computational scope}
\label{sec:problem-computation}

For fixed $k_{\max}$, the core BayesBreak workflow consists of:
\begin{enumerate}[label=(\alph*), leftmargin=2.3em]
\item precomputing all block evidences and required moment numerators;
\item running sum-product DP recursions for evidences and posterior marginals; and
\item optionally running a max-sum/backtracking recursion for a joint MAP segmentation.
\end{enumerate}
When block evaluation is constant time after prefix-sum preprocessing, this yields
$\mathcal{O}(n^2)$ block precomputation and $\mathcal{O}(k_{\max}n^2)$ DP time. These worst-case
costs make the exact method most natural for moderate sequence lengths (hundreds to low
thousands), where exact posterior quantities are still computationally feasible and often worth
the cost.

\subsection{Prediction as a first-class inferential target}
\label{sec:problem-prediction}

Beyond posterior inference on the training sequence(s), BayesBreak treats \emph{prediction for new
data} as an additional inferential target at the same level as evidence, posteriors, and point
estimators. Given a fitted model $\mathcal{M}_g$ (one per group) and new observations
$(X^{\mathrm{new}},y^{\mathrm{new}})$, the relevant quantities are (i) group-conditional
posterior-predictive likelihoods $p(y^{\mathrm{new}}\mid \mathcal{M}_g)$, (ii) group membership
posteriors $p(g\mid y^{\mathrm{new}})$, and (iii) signal predictions
$\widehat f^{\mathrm{MAP}}_g(x^\star)$ or $\widehat f^{\mathrm{Bayes}}_g(x^\star)$ at query
points $x^\star$. Section~\ref{sec:prediction} develops the prediction layer in detail, covering pointwise,
set-valued, and vector-valued new data, and separates exported-segmentation scoring from optional
full resegmentation of new data.

\section{Exponential-family segment likelihood with exposure and power weights}
\label{sec:setup}
Let $d_t$ collect family-specific known information, such as Poisson exposure, Binomial trial count,
or Gaussian precision. Let $a_t\ge 0$ be an optional power or reliability weight that, when used,
multiplies the complete log contribution. The canonical local model is
\begin{equation}
 p(y_t\mid\theta,d_t)^{a_t}
 =h(y_t,d_t)^{a_t}\exp\!\left\{a_t S(y_t,d_t)^\top\eta(\theta)-a_t E(d_t)A(\theta)\right\}.
\label{eq:EFbasic}
\end{equation}
Thus
\[
 S_{ij}=\sum_{t=i+1}^j a_tS(y_t,d_t),\qquad
 W_{ij}=\sum_{t=i+1}^j a_tE(d_t),\qquad
 H_{ij}=\sum_{t=i+1}^j a_t\log h(y_t,d_t).
\]
This notation is deliberately more general than a scalar ``weight'': a Binomial trial count changes
both the statistic and the log-partition coefficient, while a power weight scales the entire
observation contribution.

\begin{table}[H]
\centering\small
\begin{tabularx}{\textwidth}{@{}lXXX@{}}
\toprule
Family & $S(y,d)$ & $E(d)$ & Role of $d$\\ \midrule
Gaussian, known precision $\omega$ & $\omega y$ & $\omega$ & precision/exposure in the quadratic family\\
Poisson, exposure $e$ & $y$ & $e$ & mean is $e\lambda$; base measure contains $e^y/y!$\\
Binomial, trials $m$ & $y$ & $m$ & log-partition is multiplied by the trial count\\
Unweighted point observation & $T(y)$ & $1$ & $d$ is empty and $a=1$\\
\bottomrule
\end{tabularx}
\caption{Valid specializations of the exposure-aware segment likelihood. Optional $a_t$ is a separate likelihood-power weight.}
\label{tab:exposure-model}
\end{table}

\section{Ordered partitions and factorized priors}
\begin{assumption}[Factorizing ordered-partition prior]
\label{ass:factorizing-prior}
For fixed $k$ and candidate-boundary coordinates $x$, the prior on
$t=(0=t_0<\cdots<t_k=n)$ has the form
\begin{equation}
 p(t\mid k,x)=C_k(x)^{-1}
 \prod_{q=1}^{k}c_x(t_{q-1},t_q)
 \prod_{q=1}^{k-1}h_x(t_q),
\label{eq:partition-prior-factor}
\end{equation}
where $c_x(i,j)\ge0$ is a segment cohesion and $h_x(j)\ge0$ is an interior-boundary hazard.
The terminal endpoint $n$ receives no hazard factor. Define
$\gamma_x(i,j)=c_x(i,j)h_x(j)^{\mathbf1\{j<n\}}$.
\end{assumption}
The normalizer is
\begin{equation}
 C_k(x)=\sum_{t\in\mathcal T^+_{k,n}}
 \prod_{q=1}^{k}c_x(t_{q-1},t_q)\prod_{q<k}h_x(t_q),
\label{eq:Ck-setup}
\end{equation}
and the prior-adjusted local arrays are
\begin{equation}
 \widetilde A^{(r)}_{ij}=A^{(r)}_{ij}\gamma_x(i,j).
\label{eq:Atilde-setup}
\end{equation}
This factorization subsumes an index-uniform prior, a length cohesion, a local boundary-intensity
model, or their product without confusing their interpretations.
